\chapter{Preface}

The goal of this Fortran tutorial is to give a quick introduction to the most common features of the Fortran 77 programming language. A companion tutorial introduces the key enhancements found in Fortran 90. It is \emph{not} a complete reference! Many details have been omitted. The presentation focuses on scientific computations, mainly linear algebra. The outline of this tutorial was inspired by the book \emph{"Handbook for Matrix Computations"} by T.F. Coleman and C. Van Loan, published by \href{http://www.siam.org/}{SIAM}. Sections of this tutorial pertaining to Fortran 90 were based on the book \emph{"FORTRAN 77 for Engineers and Scientists with an Introduction to Fortran 90"} by L. Nyhoff and S. Leestma published by Prentice Hall.

This tutorial was designed by Erik Boman to be used in the course \emph{SCCM-001-F: Introduction to Fortran} taught at Stanford University, Winter quarter 1996. It has been modified by Sarah T. Whitlock and Paul H. Hargrove for use in the Fortran courses which have been offered under different course numbers each subsequent year. Permission to use this tutorial for educational and other non-commercial purposes is granted provided all author and copyright information is retained.

\begin{description}
    \item[Paul H. Hargrove, Stanford, December 1997.]
    \item[Sarah T. Whitlock, Stanford, January 1997.]
    \item[Erik Boman, Stanford, December 1995.]
\end{description}

\emph{Copyright \textcopyright{} 1995-7 by Stanford University. All rights reserved.}

\section*{A note on this edition}
\addcontentsline{toc}{section}{A note on this edition}

This is an adapted and extended edition of the Fortran 77 tutorial written at Stanford University by Erik Boman, Sarah T. Whitlock and Paul H. Hargrove, whose original preface appears above, unchanged. The original tutorial is still available at \url{https://web.stanford.edu/class/me200c/tutorial_77/}. This edition was prepared by Roberto Augusto G\'omez Loenzo in 2026 and is distributed free of charge, for educational and other non-commercial purposes, under the permission granted in the original preface.

The main changes with respect to the original are:

\begin{itemize}
    \item corrections of technical errors, and updates of the parts that had become outdated (compilers, tools, the status of obsolescent and deleted features in the later Fortran standards, and the resources on the Web);
    \item a new chapter on compiling with gfortran, and new or extended sections in several chapters (for example statement functions, \texttt{save}, internal files, implied do-loops, error handling in file I/O, and debugging options);
    \item exercises at the end of every chapter (the original tutorial had none);
    \item the directory \texttt{code} that accompanies this edition, with all the example programs completed so that they can be compiled and run, the tested solutions of the exercises, and a script that compiles and runs all of them.
\end{itemize}

The references above to a companion Fortran 90 tutorial and to sections on Fortran 90 belong to the original edition; they are not part of this one.

This edition is \emph{not} an official version of the Stanford tutorial, and it has not been reviewed or endorsed by its authors or by Stanford University. Any errors introduced in it are the responsibility of the adapter.

The parts taken from the original tutorial remain Copyright \textcopyright{} \mbox{1995-7} by Stanford University, under the terms of its original permission. The new material of this edition (new and rewritten text, exercises and their solutions, and the programs in the \texttt{code} directory) is Copyright \textcopyright{} 2026 Roberto Augusto G\'omez Loenzo and is licensed under the Creative Commons Attribution-NonCommercial-ShareAlike 4.0 International License (CC BY-NC-SA 4.0), \url{https://creativecommons.org/licenses/by-nc-sa/4.0/}. In short, you may copy and adapt it for non-commercial purposes, as long as you credit both the original authors and the adapter, and distribute your adaptations under the same terms.
